\documentclass[11pt]{article}
        \usepackage[margin=0.72in]{geometry}
        \usepackage[T1]{fontenc}
        \usepackage[utf8]{inputenc}
        \usepackage{lmodern}
        \usepackage{microtype}
        \usepackage{graphicx}
        \usepackage{booktabs}
        \usepackage{tabularx}
        \usepackage{array}
        \usepackage{xcolor}
        \usepackage{float}
        \usepackage{caption}
        \usepackage{enumitem}
        \usepackage{fancyhdr}
        \usepackage[numbers,sort&compress]{natbib}
        \usepackage[colorlinks=true,linkcolor=blue!55!black,citecolor=blue!55!black,urlcolor=blue!55!black]{hyperref}

        \definecolor{navy}{HTML}{163A5F}
        \definecolor{pale}{HTML}{EEF4F9}
        \definecolor{warning}{HTML}{FFF3CD}
        \definecolor{rulegray}{HTML}{B8C2CC}
        \newcommand{\auditbox}[1]{%
          \par\smallskip\noindent
          \fcolorbox{navy}{pale}{\parbox{0.955\linewidth}{\textbf{Detector audit.} #1}}%
          \par\smallskip
        }
        \newcommand{\cautionbox}[1]{%
          \par\smallskip\noindent
          \fcolorbox{rulegray}{warning}{\parbox{0.955\linewidth}{\textbf{Critical limitation.} #1}}%
          \par\smallskip
        }

        \pagestyle{fancy}
        \fancyhf{}
        \lhead{MIT--BIH record 113 beat-type atlas}
        \rhead{RR--HRV front end}
        \cfoot{\thepage}
        \setlength{\headheight}{14pt}
        \setlength{\parindent}{0pt}
        \setlength{\parskip}{4.5pt}
        \setlist[itemize]{leftmargin=1.2em,itemsep=2pt,topsep=3pt}
        \captionsetup{font=small,labelfont=bf}

        \title{\textbf{MIT--BIH Record 113: Detailed Beat-Type Atlas}\\[0.35em]
        \large Expert labels, physiological interpretation, and current QRS-detection audit}
        \author{ECG-only seizure detector project}
        \date{4 August 2026}

        \begin{document}
        \maketitle

        \begin{center}
        \fcolorbox{navy}{pale}{\parbox{0.91\textwidth}{
        \textbf{The central distinction.} The symbols in this report are expert WFDB beat annotations.
        The current RR--HRV front end detects QRS/R events and measures detector agreement; it does
        \emph{not} predict these beat classes. Therefore ``QRS detected'' is never rewritten as
        ``beat type classified correctly,'' and detector disagreement is never automatically called artifact.
        }}
        \end{center}

        \section{Record, data, and selection rule}

        The MIT--BIH Arrhythmia Database contains 48 two-lead ambulatory ECG recordings with expert-reviewed
        beat annotations \cite{goldberger2000physiobank,moody2001mitbih,physionetmitdb}. Record 113 is
        a 30-minute, 360-Hz recording from a 24-year-old woman with MLII and V1. The official notes describe predominantly normal sinus rhythm with rate variation possibly caused by a wandering atrial pacemaker, and no special noise warning. The figures use the upper MLII lead.

        The beat labels audited here are \texttt{N} ($n=1,789$), \texttt{a} ($n=6$). For each label, the displayed target is the chronologically
        middle occurrence after excluding only the 1.5-second file edges. This deterministic rule does not
        select the best-looking or best-detected beat. The comparison comes from another MIT--BIH record with
        the same expert symbol, so the report shows real published variability rather than an invented ideal.

        \subsection{Current detector path and marker meanings}

        The primary event detector is the Khamis/UNSW QRS detector; the NeuroKit gradient detector supplies
        secondary context \cite{khamis2016telehealth,makowski2021neurokit,kristof2024qrs}. The selected
        output is the original-polarity UNSW event refined to the highest positive amplitude within $\pm50$~ms.
        The inverted-polarity result is shown as context only and cannot silently replace the selected timestamp.
        A primary event is called supported when exactly one secondary event lies within $\pm50$~ms. Expert
        matching uses a separate one-to-one $\pm75$-ms audit tolerance.

        In each target panel: the red-edged diamond is the selected expert annotation; gray diamonds and labels
        are neighboring expert beats; the teal triangle is the initial UNSW event; the orange cross is the
        selected refined R timestamp; the purple plus-shaped marker is inverted-path context; and black baseline
        ticks are NeuroKit-gradient context. None of those detector markers is a beat-class prediction.

        \cautionbox{The $\pm75$-ms audit tolerance defines this project's matching experiment; it is not an
        industry declaration that any timestamp inside the window is an exact R-wave fiducial. A wide tolerance
        can hide timing jitter that matters to RR and HRV. Exact raw timing errors are therefore retained.}

        \clearpage
        \section{Full-record QRS-detection audit by expert label}

        \begin{table}[H]
        \centering
        \caption{Selected refined R timestamps versus all expert annotations of each requested label.}
        \small
        \begin{tabularx}{\textwidth}{lXrrrr}
        \toprule
        Symbol & Official meaning & Expert $n$ & Matched/$n$ & Sensitivity & 95\% Wilson interval \\
        \midrule
        \texttt{N} & Expert reference label for a normal beat. & 1,789 & 1,788/1,789 & 99.94\% & 99.68--99.99\% \\
\texttt{a} & Expert reference label for an aberrated atrial premature beat. & 6 & 6/6 & 100.00\% & 60.97--100.00\% \\
        \bottomrule
        \end{tabularx}
        \end{table}

        \begin{figure}[H]
        \centering
        \includegraphics[width=0.93\textwidth]{figures/record_113_symbol_sensitivity.png}
        \caption{The interval width, not just the point estimate, exposes how little rare labels establish.
        Intervals use the Wilson score method \cite{wilson1927probable}.}
        \end{figure}

        The full-record initial UNSW result was TP=1,794, FP=1,
        FN=1, and $F_1=0.99944$. After the current positive-amplitude
        fiducial refinement it was TP=1,794, FP=1,
        FN=1, and $F_1=0.99944$. Per-label sensitivity allocates
        matched expert beats to their symbols; false-positive detector events have no expert symbol and therefore
        cannot be represented by sensitivity alone. This table is not a beat-classification confusion matrix.

        \subsection{Exact selected-example audit}
        \begin{table}[H]
        \centering
        \caption{Timing results for the deterministic displayed examples. Errors are detector minus expert.}
        \resizebox{\textwidth}{!}{%
        \begin{tabular}{lrrrrrrrrl}
        \toprule
        Expert & Sample & Time (s) & Initial error (ms) & Selected R error (ms) & Supported? & Inverted R error (ms) & QRS detected? & Beat class output \\
        \midrule
        \texttt{N} & 331,431 & 920.642 & $+5.6$ & $0.0$ & Yes & $+16.7$ & Yes & Not produced \\
\texttt{a} & 146,714 & 407.539 & $+5.6$ & $0.0$ & Yes & $+36.1$ & Yes & Not produced \\
        \bottomrule
        \end{tabular}}
        \end{table}

        \subsection{Record-specific critical reading}
        \begin{itemize}
        \item The selected output matched all six expert a beats, but the 95\% Wilson interval remains 60.97--100.00\%. Six observations do not establish robust aberrated-atrial-beat generalization.
\item The official record has no special noise warning. Successful QRS localization here therefore tests clean morphology more than severe artifact tolerance.
\item Both displayed selected refined timestamps are exact at the 360-Hz sample grid, while the inverted path is later. This supports the selected polarity for these examples only, not an automatic global routing rule.
        \end{itemize}

        \clearpage
\section{\texttt{N}: normal beat}

\begin{figure}[H]
\centering
\includegraphics[width=\textwidth]{figures/record_113_N_comparison.png}
\caption{Left: the chronologically middle expert \texttt{N} occurrence in record 113
(sample 331,431, 920.642~s), with current detector outputs. Right: an independent
published MIT--BIH record 100 occurrence carrying the same expert symbol. Neighboring
reference symbols are shown to preserve rhythm context \cite{moody2001mitbih,physionetmitdb}.}
\end{figure}

\textbf{Official symbol meaning.} Expert reference label for a normal beat.

\textbf{Why this type occurs and what tends to shape the waveform.} A sinus impulse normally depolarizes the atria first, producing the P wave, then reaches the ventricles through the AV node and His--Purkinje system, producing the QRS; ventricular repolarization produces the T wave \cite{kligfield2007standardization,surawicz2009conduction}. The database symbol describes the expert beat class. It does not require the largest QRS deflection to point upward in every lead.

\textbf{How to read the two strips.} Polarity and amplitude depend on the projection of the cardiac electrical vector onto the selected lead. A negative or biphasic complex may therefore still carry an expert normal-beat label. The comparison strip demonstrates allowed appearance variation; it is not a universal template. The right-hand strip is not a
synthetic textbook ideal and was not selected as a detector success; it is another expert-labeled
observation from the published database. Similar labels can legitimately have different waveforms.

\auditbox{For the deterministic middle N example, the UNSW initial event was $+5.6$~ms from the expert annotation; the selected original-polarity refined timestamp was $0.0$~ms away and was within the predefined $\pm75$-ms audit tolerance. Exact NeuroKit-gradient support within $\pm50$~ms of the primary QRS event was yes. The inverted-context refined timestamp was $+16.7$~ms away. Across the full record, the selected refined output matched 1788 of 1789 expert N annotations. This is QRS localization, not beat-type classification.}

\cautionbox{A single-lead waveform cannot establish that the entire 12-lead ECG is normal, and QRS detection cannot establish the N class.}
\clearpage
\section{\texttt{a}: aberrated atrial premature beat}

\begin{figure}[H]
\centering
\includegraphics[width=\textwidth]{figures/record_113_a_comparison.png}
\caption{Left: the chronologically middle expert \texttt{a} occurrence in record 113
(sample 146,714, 407.539~s), with current detector outputs. Right: an independent
published MIT--BIH record 201 occurrence carrying the same expert symbol. Neighboring
reference symbols are shown to preserve rhythm context \cite{moody2001mitbih,physionetmitdb}.}
\end{figure}

\textbf{Official symbol meaning.} Expert reference label for an aberrated atrial premature beat.

\textbf{Why this type occurs and what tends to shape the waveform.} The impulse is premature and atrial, but part of the ventricular conduction system is still refractory. Ventricular activation therefore follows a bundle-branch or fascicular-block pattern, widening or changing the QRS even though the origin is supraventricular \cite{samesima2022ecg,escudero2021wide}.

\textbf{How to read the two strips.} This label illustrates why `wide' does not automatically mean ventricular. The timing of the early atrial impulse, any altered P wave, and the relationship to neighboring cycles help distinguish aberrancy from a PVC. The right-hand strip is not a
synthetic textbook ideal and was not selected as a detector success; it is another expert-labeled
observation from the published database. Similar labels can legitimately have different waveforms.

\auditbox{For the deterministic middle a example, the UNSW initial event was $+5.6$~ms from the expert annotation; the selected original-polarity refined timestamp was $0.0$~ms away and was within the predefined $\pm75$-ms audit tolerance. Exact NeuroKit-gradient support within $\pm50$~ms of the primary QRS event was yes. The inverted-context refined timestamp was $+36.1$~ms away. Across the full record, the selected refined output matched 6 of 6 expert a annotations. This is QRS localization, not beat-type classification.}

\cautionbox{There are only six such annotations in record 113. Detection of their QRS complexes is not evidence that the architecture can distinguish a from V. The local denominator is only 6; always report the raw 6/6 rather than presenting the percentage alone.}

        \clearpage
        \section{Conclusions for record 113}
        \begin{enumerate}
        \item This report validates QRS-event localization against expert timestamps. It does not validate
        normal, atrial, ventricular, conduction-pattern, or junctional classification;
        seizure detection; or artifact detection.
        \item The full raw denominators and Wilson intervals must accompany percentages. A 100\% result with
        one or two beats remains weak evidence.
        \item The independent comparison strips show that one expert label can have multiple lead-dependent
        shapes. Template resemblance alone is not ground truth.
        \item Original and inverted refinements are both displayed because morphology and polarity can move the
        chosen positive maximum. The inverted path remains context until a routing rule is prospectively specified
        and independently validated.
        \item Detector support is a reliability measurement, not a veto and not an artifact label. Expert
        annotations remain the evaluation reference.
        \end{enumerate}

        \bibliographystyle{unsrtnat}
        \bibliography{MITBIH_Beat_Type_Atlases}
        \end{document}
